\documentclass[11pt,a4paper]{article}

\usepackage[utf8]{inputenc}

\usepackage[T1]{fontenc}

\usepackage[polish]{babel}

\usepackage{lmodern}

\usepackage{geometry}

\geometry{margin=27mm}

\usepackage[hidelinks]{hyperref}

\title{Wnioskowanie o czasie recovery przy cenzorowaniu i niepełnej obserwacji}

\author{Mateusz Pomianek}

\date{Szkic naukowy, wersja 1.0.0 | 28 listopada 2026}

\begin{document}

\maketitle

\begin{abstract}

Proponowana miara główna to bezwzględne obciążenie oszacowania skumulowanej częstości poprawnego powrotu do 30 s.

\end{abstract}

\section{Wprowadzenie}

Zdarzenie konkurencyjne zmienia interpretację prawdopodobieństwa zdarzenia docelowego (Austin et al., 2016). Wynik Briera pozwala oceniać jakość prognoz przeżycia przy odpowiedniej korekcie cenzorowania (Graf et al., 1999).

\section{Cel i pytania badawcze}

Hipotezy robocze wymagają zatwierdzenia przed rozpoczęciem właściwej weryfikacji.

\paragraph{H1.} Przy dodatniej częstości definitywnego porzucenia zadania estymator skumulowanej częstości powrotu uwzględniający porzucenie ma mniejsze bezwzględne obciążenie prawdopodobieństwa powrotu do 30 s niż dopełnienie estymatora Kaplana i Meiera cenzorujące porzucenie.

\section{Status koncepcji i granice dowodów}

Koncepcja ma status IDEA i DEFER. Nie wskazano zatwierdzonego protokołu, kanonicznego projektu źródłowego ani zbioru obserwacji. Właścicielem naukowym szkicu jest Mateusz Pomianek. Wskazanie to nie potwierdza przeprowadzenia badania ani nie rozstrzyga autorstwa przyszłej publikacji.

\section{Materiał i metody}

Pierwszorzędowy wskaźnik: bezwzględne obciążenie oszacowania skumulowanej częstości poprawnego powrotu do 30 s.

W symulacji referencją jest prawdopodobieństwo zadane przez generator. Należy oddzielić administracyjny koniec okna, definitywne porzucenie i utratę pomiaru; zależną od niewidocznego stanu utratę obserwacji zbada analiza wrażliwości. Porównanie obejmie scenariusz kontrolny bez zdarzeń konkurencyjnych.

Przed badaniem z udziałem ludzi należy zatwierdzić wymagania etyczne, zasady zgody i ochronę danych.

\section{Plan analizy}

Zamrożenie estymandu, reguł oceny jakości i obsługi braków nastąpi przed dostępem do danych wynikowych. Raport obejmie oszacowanie kontrastu głównego z przedziałem niepewności i pełne mianowniki. Jeśli przewidziano hipotezy drugorzędowe, przed analizą należy ustalić sposób kontroli wielokrotności.

\section{Wyniki}

W niniejszym pliku nie umieszczono wyników empirycznych, statystyk ani nieprzeprowadzonych analiz. Sekcja wymaga uzupełnienia po uzyskaniu i weryfikacji materiału dowodowego.

\section{Dyskusja i ograniczenia}

Ograniczenia obejmują możliwe błędy pomiaru, dobór próby, alternatywne mechanizmy i ograniczoną możliwość uogólniania. Korelacja nie dowodzi przyczynowości; poprawność techniczna nie jest dowodem znaczenia klinicznego.

\section{Wnioski}

Hipotezy i sposób pomiaru stanowią podstawę oceny wykonalności oraz przygotowania właściwego protokołu.

\section*{Bibliografia}

{\setlength{\parindent}{0pt}
\setlength{\parskip}{5pt}
\hangindent=1.27cm
Austin, P. C., Lee, D. S., \& Fine, J. P. (2016). Introduction to the analysis of survival data in the presence of competing risks. \textit{Circulation, 133}(6), 601--609. \url{https://doi.org/10.1161/CIRCULATIONAHA.115.017719}\par
\hangindent=1.27cm
Graf, E., Schmoor, C., Sauerbrei, W., \& Schumacher, M. (1999). Assessment and comparison of prognostic classification schemes for survival data. \textit{Statistics in Medicine, 18}(17--18), 2529--2545. \url{https://pubmed.ncbi.nlm.nih.gov/10474158/}\par
}

\end{document}
